\documentclass[11pt]{article}
\usepackage{amsmath,amssymb,amsthm}
\usepackage{booktabs}
\usepackage{geometry}
\geometry{margin=1in}

\newcommand{\R}{\mathbb{R}}
\newcommand{\ip}[2]{\langle #1, #2 \rangle}
\newcommand{\half}{\tfrac{1}{2}}
\DeclareMathOperator{\Exp}{Exp}

\newtheorem{definition}{Definition}
\newtheorem{theorem}{Theorem}
\newtheorem{proposition}{Proposition}
\newtheorem{corollary}[theorem]{Corollary}
\newtheorem{lemma}{Lemma}

\title{High-Order Geodesic Integration on Hessian Manifolds\\
Without Third Derivatives:\\
A Tutorial}
\author{}
\date{}

\begin{document}
\maketitle

\begin{abstract}
We describe a method for integrating geodesics on Hessian manifolds
(the Riemannian geometry induced by the Hessian of a convex function)
that achieves fourth-order accuracy using only gradient and Hessian
evaluations---no third derivatives (Christoffel symbols) are needed.
The method composes a second-order ``Legendre midpoint'' step using
the classical Yoshida construction.  We explain the ideas using only
the geodesic equation and the Legendre transform, without the
machinery of affine connections.
\end{abstract}

\section{Setup}

Let $F : \Omega \to \R$ be a strictly convex function on an open
convex domain $\Omega \subset \R^n$.  The \emph{Hessian metric} is
the Riemannian metric $g(s) = \nabla^2 F(s)$, which is positive
definite by strict convexity.

The \emph{geodesic} is the curve $\gamma(t)$ that locally minimizes
the length
\[
  \ell(\gamma) = \int_0^1 \sqrt{\dot\gamma(t)^\top \nabla^2 F(\gamma(t)) \, \dot\gamma(t)} \; dt.
\]
The Euler--Lagrange equation for this variational problem gives the
geodesic ODE:
\begin{equation}\label{eq:geodesic}
  \ddot\gamma^i + \frac{1}{2} \bigl[\nabla^2 F(\gamma)\bigr]^{-1}_{ik}
  \sum_{j,\ell} F_{j\ell k}(\gamma) \, \dot\gamma^j \dot\gamma^\ell = 0,
\end{equation}
where $F_{j\ell k} = \partial^3 F / \partial s_j \partial s_\ell \partial s_k$.
Integrating this ODE requires the \emph{third derivatives} of $F$
(the ``Christoffel symbols'' of the metric).

\paragraph{Goal.}
We want to integrate~\eqref{eq:geodesic} \emph{without} computing
$F_{j\ell k}$.  We will use only:
\begin{itemize}
  \item $\nabla F(s)$ (gradient), and
  \item $\nabla^2 F(s)$ (Hessian).
\end{itemize}

\section{The Legendre transform and dual coordinates}

The key tool is the \emph{Legendre transform}.  Define the dual
coordinate
\[
  \lambda = -\nabla F(s).
\]
This is a diffeomorphism from $\Omega$ to a dual domain $\Omega^*$
(the image of $-\nabla F$).  Its inverse, which we write
$s = \mathcal{L}(\lambda)$, satisfies $-\nabla F(\mathcal{L}(\lambda)) = \lambda$.

\paragraph{Two coordinate systems.}
The \emph{same} Riemannian manifold $(\Omega, \nabla^2 F)$ can be
described in two coordinate systems:
\begin{itemize}
  \item \textbf{Primal coordinates} $s$: the original variables.
  \item \textbf{Dual coordinates} $\lambda = -\nabla F(s)$: the
    Legendre-transformed variables.
\end{itemize}
A straight line in primal coordinates ($s(t) = s_0 + tv$) is
\emph{not} a geodesic.  A straight line in dual coordinates
($\lambda(t) = \lambda_0 + tw$) is also not a geodesic.  But
the true geodesic lies ``between'' these two approximations in a
precise sense.

\section{The Legendre midpoint step}

The Legendre midpoint is a second-order approximation to the geodesic
that uses one primal step and one dual step:

\begin{enumerate}
  \item \textbf{Primal half-step.}  Move halfway along a straight
    line in primal coordinates:
    \[
      s_{1/2} = s_0 + \frac{h}{2} v.
    \]

  \item \textbf{Evaluate dual at midpoint.}  Compute the dual
    coordinate at $s_{1/2}$:
    \[
      \lambda_{1/2} = -\nabla F(s_{1/2}).
    \]
    Also let $\lambda_0 = -\nabla F(s_0)$.

  \item \textbf{Dual extrapolation.}  Take a full step in dual
    coordinates, using the midpoint velocity:
    \[
      \lambda_1 = 2\lambda_{1/2} - \lambda_0.
    \]
    (This is the midpoint rule: $\dot\lambda \approx (\lambda_{1/2} - \lambda_0)/(h/2)$,
    so $\lambda_1 = \lambda_0 + h \cdot \dot\lambda$.)

  \item \textbf{Invert the Legendre transform.}  Find $s_1$ such that
    $-\nabla F(s_1) = \lambda_1$.  This is a nonlinear solve
    (Newton's method on $\nabla F(s) + \lambda_1 = 0$, using the
    Hessian $\nabla^2 F$ as the Jacobian).
\end{enumerate}

\subsection{Second-order accuracy}
\label{sec:legendre-proof}

\begin{theorem}[Legendre midpoint matches the geodesic to second order]
\label{thm:legendre-second-order}
Let $F$ be a $C^3$ barrier on the interior of a proper cone, with
Hessian $H_0 = \nabla^2 F(s_0)$.  Let $\gamma$ be the geodesic
from $s_0$ with velocity $v$ (i.e., $\gamma(0) = s_0$,
$\dot\gamma(0) = v$, and $\gamma$ satisfies the geodesic
equation~\eqref{eq:geodesic}).  Let $s_1$ be the output of the
Legendre midpoint step with step size $h$.  Then
\[
  s_1 = \gamma(h) + O(h^3).
\]
More precisely: the Taylor expansions of $s_1$ and $\gamma(h)$
in powers of $h$ agree through order $h^2$, with
\[
  s_1 = s_0 + hv + \frac{h^2}{2}\ddot\gamma(0) + O(h^3),
\]
where $\ddot\gamma(0) = -\frac{1}{2}H_0^{-1}T(v,v)$ and
$T(v,v)_i = \sum_{jk} F_{ijk}(s_0)\, v_j v_k$.
\end{theorem}

\begin{proof}
Write $H = H_0$, $\lambda_0 = -\nabla F(s_0)$, and
$T(v) = T(v,v)$ for brevity.  We compute the Taylor expansion
of the Legendre midpoint output $s_1$ and show its $O(h^2)$ coefficient
equals $\ddot\gamma(0)$.

Write $s_1 = s_0 + hv + \frac{h^2}{2}a + O(h^3)$ with unknown
acceleration $a$.  The proof determines $a$ by expanding each
step of the algorithm.

\medskip\noindent
\textit{Step 1 (primal half-step).}\quad
$s_{1/2} = s_0 + \frac{h}{2}v$.

\medskip\noindent
\textit{Step 2 (dual at midpoint).}\quad
Taylor-expand $\nabla F$ at $s_{1/2}$:
\[
  \lambda_{1/2} = -\nabla F\!\left(s_0 + \tfrac{h}{2}v\right)
  = \lambda_0 - \tfrac{h}{2}Hv - \tfrac{h^2}{8}T(v) + O(h^3).
\]
The third-derivative term $T(v)$ appears from the nonlinearity
of $\nabla F$---this is where curvature enters.

\medskip\noindent
\textit{Step 3 (dual extrapolation).}\quad
\[
  \lambda_1 = 2\lambda_{1/2} - \lambda_0
  = \lambda_0 - hHv - \tfrac{h^2}{4}T(v) + O(h^3).
\]

\medskip\noindent
\textit{Step 4 (inversion).}\quad
The algorithm defines $s_1$ by $-\nabla F(s_1) = \lambda_1$.
We do not solve this equation; instead we determine the $O(h^2)$
coefficient $a$ by showing consistency.

Write $u = s_1 - s_0 = hv + \frac{h^2}{2}a + O(h^3)$ and
Taylor-expand $\nabla F(s_0 + u)$ around~$s_0$:
\[
  \nabla F(s_0 + u)
  = \nabla F(s_0) + Hu + \tfrac{1}{2}D^2(\nabla F)[u,u] + O(\|u\|^3).
\]
The linear term gives $Hu = hHv + \frac{h^2}{2}Ha$.
The quadratic term is $\frac{1}{2}T(u,u)$, where
$T(u,u)_i = \sum_{jk}F_{ijk}u_j u_k$.  Since $u = hv + O(h^2)$,
we have $T(u,u) = h^2 T(v,v) + O(h^3) = h^2 T(v) + O(h^3)$.
Combining (and using $\lambda_0 = -\nabla F(s_0)$):
\[
  -\nabla F(s_1) = \lambda_0 - hHv
    - \tfrac{h^2}{2}Ha - \tfrac{h^2}{2}T(v) + O(h^3).
\]

\medskip\noindent
\textit{Matching.}\quad
Setting $-\nabla F(s_1) = \lambda_1$ and comparing the $O(h^2)$
coefficients from Step~3 and Step~4:
\[
  \underbrace{-\tfrac{1}{2}Ha - \tfrac{1}{2}T(v)}_{\text{from }
    -\nabla F(s_1)}
  = \underbrace{-\tfrac{1}{4}T(v)}_{\text{from }\lambda_1}.
\]
Solving: $Ha = -\frac{1}{2}T(v)$, i.e.,
$a = -\frac{1}{2}H^{-1}T(v) = \ddot\gamma(0)$. \qedhere
\end{proof}

\paragraph{Where does the curvature come from?}
The method never computes $F_{ijk}$ (third derivatives) explicitly.
The curvature information is encoded in the
\emph{nonlinearity of $\nabla F$}: the Taylor expansion of
$\nabla F(s_0 + \frac{h}{2}v)$ contains the $\frac{h^2}{8}T(v)$ term
automatically.  The midpoint extrapolation (Step~3) amplifies it to
$\frac{h^2}{4}T(v)$, and the inversion (Step~4) contributes another
$\frac{h^2}{2}T(v)$.  Together:
$-\frac{1}{4} + \frac{1}{2} = \frac{1}{4}$, which after the $H^{-1}$
solve gives $\frac{1}{2}H^{-1}T$---exactly the geodesic acceleration.

\paragraph{Analogy with mechanics.}
This is the same principle as the St\"ormer--Verlet method: a half-step
in position, a full step in momentum, a half-step in position.  The
Legendre transform $\lambda = -\nabla F(s)$ plays the role of
$p = \partial L / \partial \dot{q}$.  In both cases, the second-order
accuracy comes from sampling the nonlinearity at the midpoint.

\paragraph{Cost.}
Each Legendre midpoint step requires:
\begin{itemize}
  \item 2 gradient evaluations ($\nabla F$ at $s_0$ and $s_{1/2}$),
  \item 1 Newton solve ($\sim$3--5 Hessian evaluations for the inversion).
\end{itemize}
No third derivatives.

\section{Exactness for the log barrier: Legendre midpoint = Pad\'e}

For the LP barrier $F(s) = -\sum_i \log s_i$, the Legendre midpoint
is not merely a second-order approximation---it is \emph{exactly}
the $[1/1]$ Pad\'e approximant of the geodesic exponential.

\begin{theorem}\label{thm:pade}
Let $F(s) = -\sum_i \log s_i$ (the logarithmic barrier on $\R^n_+$).
The geodesic from $s_0 > 0$ with velocity $v$ is
$\gamma_i(h) = s_{0,i} \exp(h v_i / s_{0,i})$.
The Legendre midpoint step from $s_0$ with step size $h$ and direction $v$
gives
\[
  s_{1,i} = s_{0,i} \cdot \frac{1 + d_i/2}{1 - d_i/2},
  \qquad d_i = \frac{h v_i}{s_{0,i}},
\]
which is the $[1/1]$ Pad\'e approximant of $\exp(d_i)$:
\[
  \frac{1 + d/2}{1 - d/2} = \exp(d) + O(d^3).
\]
\end{theorem}

\begin{proof}
The barrier has $\nabla F(s) = -s^{-1}$ (componentwise) and
$\nabla^2 F(s) = \operatorname{diag}(s^{-2})$.  The Legendre transform
is $\lambda = -\nabla F(s) = s^{-1}$, with inverse $s = \lambda^{-1}$.

The four Legendre midpoint steps (componentwise, dropping the index $i$):
\begin{enumerate}
  \item $s_{1/2} = s_0 + \frac{h}{2}v$.

  \item $\lambda_{1/2} = s_{1/2}^{-1} = (s_0 + \frac{h}{2}v)^{-1}$, \quad
        $\lambda_0 = s_0^{-1}$.

  \item $\lambda_1 = 2\lambda_{1/2} - \lambda_0
        = \frac{2}{s_0 + \frac{h}{2}v} - \frac{1}{s_0}
        = \frac{2s_0 - (s_0 + \frac{h}{2}v)}{s_0(s_0 + \frac{h}{2}v)}
        = \frac{s_0 - \frac{h}{2}v}{s_0(s_0 + \frac{h}{2}v)}$.

  \item $s_1 = \lambda_1^{-1}
        = \frac{s_0(s_0 + \frac{h}{2}v)}{s_0 - \frac{h}{2}v}
        = s_0 \cdot \frac{1 + d/2}{1 - d/2}$,
        \quad where $d = hv/s_0$.
\end{enumerate}
The $[1/1]$ Pad\'e approximant of $e^d$ is $(1+d/2)/(1-d/2)$, which
satisfies $(1+d/2)/(1-d/2) = 1 + d + d^2/2 + \cdots = e^d + O(d^3)$.
\end{proof}

\begin{corollary}
For $F(s) = -\log\det(S)$ on the PSD cone $\mathbb{S}^n_+$, the Legendre
midpoint reduces to the $[1/1]$ matrix Pad\'e approximant of $\exp(D)$
when $S_0$ and $V$ are simultaneously diagonalizable.
\end{corollary}
\begin{proof}
In the joint eigenbasis, $F$ decomposes into $-\sum_i \log \sigma_i$
where $\sigma_i$ are the eigenvalues.  Each component follows
Theorem~\ref{thm:pade} independently.
\end{proof}

\paragraph{Remark.}
The Yoshida fourth-order composition of three Pad\'e $[1/1]$ steps
produces a higher-order rational approximation to $\exp(d)$---a
$[2/2]$-like Pad\'e obtained by composition rather than direct
construction.  For the log barrier, the Yoshida--Legendre method is
therefore a \emph{high-order rational approximation to the geodesic
exponential} that generalizes seamlessly from symmetric cones (where
it reduces to Pad\'e) to arbitrary barriers (where it uses the
Legendre transform in place of the matrix exponential).

\section{The Yoshida fourth-order composition}

The Legendre midpoint is second-order: the error is $O(h^3)$.  To get
fourth-order accuracy ($O(h^5)$ error), we use the \emph{Yoshida
construction} (1990).

\paragraph{Idea.}
If $\Psi_h$ is a second-order method (so $\Psi_h = \Phi_h + O(h^3)$
where $\Phi_h$ is the exact flow), then the composition
\begin{equation}\label{eq:yoshida}
  \Psi_h^{[4]} = \Psi_{x_1 h} \circ \Psi_{x_0 h} \circ \Psi_{x_1 h}
\end{equation}
is fourth-order, provided
\[
  2x_1 + x_0 = 1 \quad \text{and} \quad 2x_1^3 + x_0^3 = 0.
\]
The unique real solution is
\[
  x_1 = \frac{1}{2 - 2^{1/3}} \approx 1.3512, \qquad
  x_0 = \frac{-2^{1/3}}{2 - 2^{1/3}} \approx -1.7024.
\]

\paragraph{Why does this work?}
The second-order method has error $\Psi_h = \Phi_h + c_3 h^3 + O(h^5)$
(no $h^4$ term because $\Psi$ is symmetric/time-reversible).  The
composition~\eqref{eq:yoshida} gives
\[
  \Psi_h^{[4]} = \Phi_h + c_3(2x_1^3 + x_0^3) h^3 + O(h^5).
\]
The condition $2x_1^3 + x_0^3 = 0$ kills the cubic error.  The
negative coefficient $x_0 < 0$ corresponds to a ``backward'' step
that cancels the forward steps' third-order error.

\paragraph{Applied to the Legendre midpoint.}
The Yoshida fourth-order geodesic step is:
\begin{enumerate}
  \item Legendre step with step size $x_1 \cdot h$ \quad (forward),
  \item Legendre step with step size $x_0 \cdot h$ \quad (backward correction),
  \item Legendre step with step size $x_1 \cdot h$ \quad (forward).
\end{enumerate}
Each Legendre midpoint step updates $(s, \lambda, v)$: position, dual, and
velocity.  The velocity at the end of each substep is recovered from
the position change during the second primal half-step:
$v_{\text{new}} = (s_{\text{new}} - s_{1/2}) / (h_i/2)$.

\paragraph{Cost.}
Three Legendre steps $= 3 \times (2$ gradient $+ 1$ Newton$)$
$= 6$ gradient evaluations $+ 3$ Newton solves.  For a 3D cone
(e.g., the exponential cone), this is $\sim$150 floating-point
operations.  No third derivatives at any point.

\section{Summary of integrators}

\begin{center}
\renewcommand{\arraystretch}{1.2}
\begin{tabular}{@{}llccc@{}}
\toprule
\textbf{Method} & \textbf{Order} & \textbf{Needs $\partial^3 F$} &
  \textbf{Newton solves} & \textbf{Error at $\alpha\!=\!0.1$} \\
\midrule
Euler & 1 & No & 0 & $2.3 \times 10^{-4}$ \\
Legendre midpoint & 2 & No & 1 & $7.7 \times 10^{-7}$ \\
Yoshida--Legendre & 4 & No & 3 & $2.5 \times 10^{-9}$ \\
\midrule
St\"ormer--Verlet (8 steps) & 2 & Yes & 0 & $3.3 \times 10^{-7}$ \\
Primal-dual Verlet (8 steps) & 2+ & Yes & 1 & $3.1 \times 10^{-8}$ \\
\bottomrule
\end{tabular}
\end{center}

The Yoshida--Legendre method achieves fourth-order accuracy
($2.5 \times 10^{-9}$) without any third-derivative computation,
outperforming even the primal-dual Verlet ($3.1 \times 10^{-8}$)
which uses 8 substeps with full Christoffel symbol evaluations.

\section{Why Yoshida and not Gauss--Legendre?}

The Yoshida construction composes three $[1/1]$ Pad\'e steps to get
fourth-order accuracy.  A natural question: can we instead use the
$[2/2]$ Pad\'e directly?  For the log barrier, the $[2/2]$ Pad\'e of
$\exp(d)$ has optimal error constants (about $40\times$ smaller than
Yoshida at fifth order).  The $[2/2]$ Pad\'e corresponds to the
2-stage Gauss--Legendre implicit Runge--Kutta method.

\subsection{The obstruction}

The Legendre midpoint works because the Legendre transform provides
an \emph{exact} dual evaluation at the midpoint: $\lambda_{1/2} = -\nabla F(s_{1/2})$.
From the Taylor expansion proof (Section~\ref{sec:legendre-proof}), the
curvature enters through the nonlinearity of $\nabla F$ at the midpoint.
The \emph{symmetry} of the midpoint ($c = 1/2$) is essential: it is the
unique point where the Strang splitting of primal and dual flat flows
gives second-order accuracy.

The 2-stage Gauss--Legendre method evaluates at $c_1 = \frac{1}{2} - \frac{\sqrt{3}}{6}
\approx 0.21$ and $c_2 = \frac{1}{2} + \frac{\sqrt{3}}{6} \approx 0.79$.
At these non-symmetric points, we can evaluate $\lambda_i = -\nabla F(s_i)$
(the dual position), but we cannot recover the \emph{velocity} $v_i$ at the
stage points from the gradient map alone.

\paragraph{The velocity problem.}
At the midpoint ($c = 1/2$), the dual extrapolation
$\lambda_1 = 2\lambda_{1/2} - \lambda_0$ is the \emph{midpoint rule}
for the integral $\lambda_1 = \lambda_0 + \int_0^h \dot\lambda \, dt$.
The midpoint rule is second-order accurate: the error is $O(h^3)$
because the first-order error term vanishes when the sample point is
at the center of the interval.

At a general GL node $c_i \neq 1/2$, the analogous extrapolation
$\lambda_1 = \lambda_0 + (\lambda_{c_i} - \lambda_0)/c_i$ is
\emph{not} a midpoint rule---it is a forward-Euler extrapolation
from an off-center sample point, which is only \emph{first-order}
accurate ($O(h^2)$ error).  The GL consistency conditions couple
$v_1$ and $v_2$ through the Butcher matrix $a_{ij}$, but without
the geodesic acceleration (which requires $\partial^3 F$), the
coupled system is underdetermined.

\paragraph{What goes wrong.}
We implemented the 2-stage GL with Legendre-derived velocities and
fixed-point iteration.  The result: the iteration does not converge
to the geodesic.  The error ($1.3 \times 10^{-4}$) is comparable to
Euler, indicating that the Legendre velocity approximation destroys
the fourth-order accuracy of the GL method.

\paragraph{Why Yoshida avoids this.}
The Yoshida construction sidesteps the velocity problem entirely:
it composes three \emph{midpoint} steps (where the Legendre trick is
exact) with algebraic coefficients.  The midpoint is the only node
where the Legendre transform provides a second-order velocity without
the ODE acceleration.  Yoshida is the optimal way to achieve fourth
order from this midpoint-only primitive.

\begin{center}
\renewcommand{\arraystretch}{1.2}
\begin{tabular}{@{}lccl@{}}
\toprule
\textbf{Method} & \textbf{Order} & \textbf{Needs $\partial^3 F$} & \textbf{Error at $\alpha = 0.1$} \\
\midrule
Legendre midpoint ($= [1/1]$ Pad\'e) & 2 & No & $7.7 \times 10^{-7}$ \\
Yoshida--Legendre (3 midpoint steps) & 4 & No & $2.5 \times 10^{-9}$ \\
Gauss--Legendre 2-stage (attempted) & 4 & \textbf{Yes}${}^*$ & fails \\
\midrule
$[2/2]$ Pad\'e (optimal, log barrier) & 4 & No${}^\dagger$ & $\sim\!10^{-10}$ (theory) \\
\bottomrule
\end{tabular}
\end{center}
\noindent
${}^*$\,Without $\partial^3 F$, the GL stage velocities are only first-order. \\
${}^\dagger$\,The $[2/2]$ Pad\'e is a closed-form rational function of $d$,
specific to $F = -\log$.  It does not generalize to arbitrary barriers.

\paragraph{Summary.}
The midpoint ($c = 1/2$) is special: it is the unique evaluation point
where the dual extrapolation $\lambda_1 = \lambda_0 + (\lambda_c - \lambda_0)/c$
reduces to the \emph{midpoint quadrature rule}, which is second-order
accurate.  At any other $c$, the extrapolation is first-order (forward
Euler from an off-center point), and the resulting velocity is too
inaccurate for a fourth-order method.  Yoshida achieves fourth order
by composing three midpoint evaluations---the only node where the
Legendre transform gives second-order velocity---with algebraic
weights that cancel the third-order error.

\section{Connection to Hamiltonian mechanics}

The structure exploited here is the same as in symplectic mechanics:

\begin{center}
\begin{tabular}{@{}ll@{}}
\toprule
\textbf{Hamiltonian mechanics} & \textbf{Hessian manifold} \\
\midrule
Position $q$ & Primal coordinate $s$ \\
Momentum $p = \partial L / \partial \dot{q}$ & Dual coordinate $\lambda = -\nabla F(s)$ \\
Kinetic energy $T(p)$ & Primal ``straight line'' \\
Potential energy $V(q)$ & Dual ``straight line'' \\
Verlet = split $T + V$ & Legendre = split primal + dual \\
Yoshida on Verlet & Yoshida on Legendre \\
\bottomrule
\end{tabular}
\end{center}

In mechanics, the St\"ormer--Verlet splits the Hamiltonian $H = T(p) + V(q)$
into separately integrable kinetic and potential parts.  This works
because $T$ depends only on $p$ and $V$ depends only on $q$
(``separable'').

On a Hessian manifold, the geodesic Hamiltonian
$\mathcal{H}(s, p) = \frac{1}{2} p^\top [\nabla^2 F(s)]^{-1} p$
is \emph{not} separable (the ``kinetic energy'' depends on $s$).
But the Legendre transform provides a different splitting: primal
coordinates (where straight lines are ``free'') and dual coordinates
(where straight lines are also ``free'').  The Legendre midpoint
splits along this primal-dual axis instead of the kinetic-potential
axis.

The Yoshida construction is purely algebraic---it works for any
symmetric second-order method, regardless of the underlying geometry.
Applied to the Legendre midpoint, it gives fourth-order accuracy on
any Hessian manifold, for any strictly convex $F$, without ever
computing the curvature.

\section{Error bounds from self-concordance}

Self-concordance of the barrier provides \emph{a priori} error bounds
for all the integrators, in terms of a single scalar parameter
$\delta = h \|v\|_H$ (the step size in the local Hessian norm).

\subsection{Self-concordance inequality}

A $\nu$-logarithmically homogeneous self-concordant barrier $F$
satisfies:
\begin{equation}\label{eq:sc}
  |D^3 F(s)[u, v, w]| \leq 2\, \|u\|_H \|v\|_H \|w\|_H
\end{equation}
for all $u, v, w$, where $\|u\|_H = \sqrt{u^\top \nabla^2 F(s)\, u}$
is the local Hessian norm.  This is the polarized form of
$|D^3 F[v,v,v]| \leq 2\|v\|_H^3$.

\subsection{Tangent line (1st order)}

The tangent line $L(h) = s_0 + hv$ approximates the geodesic
$\gamma(h) = s_0 + hv + \frac{h^2}{2}\ddot{s} + O(h^3)$.
The leading error is
\[
  \gamma(h) - L(h) = -\frac{h^2}{4} H^{-1} T(v,v) + O(h^3),
\]
where $T(v,v)_i = \sum_{jk} F_{ijk} v_j v_k$.

\begin{theorem}[Tangent line error bound]\label{thm:euler-bound}
Let $F$ be a $\nu$-LHSCB with self-concordance
parameter $M_f$ (so~\eqref{eq:sc} becomes $|D^3 F[u,v,w]| \leq
2M_f\|u\|_H\|v\|_H\|w\|_H$).  Let $\gamma$ be the geodesic
from $s_0$ with velocity $v$, and set $\delta = h\|v\|_{H_0}$,
$H_0 = \nabla^2 F(s_0)$.  If $M_f \delta < \frac{1}{4}$, then
\begin{equation}\label{eq:tangent-bound}
  \|\gamma(h) - (s_0 + hv)\|_{H_0} \leq M_f\,\delta^2.
\end{equation}
\end{theorem}

\begin{proof}
\textit{Step 1: geodesic acceleration.}\quad
The geodesic equation gives $\ddot\gamma = -\frac{1}{2}H(\gamma)^{-1}
T(\dot\gamma,\dot\gamma)$, so $\|\ddot\gamma\|_{H(\gamma)}
= \frac{1}{2}\|T(\dot\gamma,\dot\gamma)\|_{H(\gamma)^{-1}}$.
For any $w$ with $\|w\|_{H(\gamma)} = 1$:
\[
  |w^\top T(\dot\gamma,\dot\gamma)|
  = |D^3F[\dot\gamma,\dot\gamma,w]|
  \leq 2M_f\|\dot\gamma\|_{H(\gamma)}^2.
\]
Since geodesics have constant speed
$\|\dot\gamma(t)\|_{H(\gamma(t))} = \|v\|_{H_0}$, we obtain
\begin{equation}\label{eq:accel-inst}
  \|\ddot\gamma(t)\|_{H(\gamma(t))} \leq M_f\|v\|_{H_0}^2
\end{equation}
for all $t$ in the domain of $\gamma$.

\medskip\noindent
\textit{Step 2: trajectory containment.}\quad
Set $\beta(t) = M_f\|\gamma(t) - s_0\|_{H_0}$.  On any interval
where $\beta < \frac{1}{2}$, the Hessian comparison theorem gives
$(1-\beta)^2 H_0 \preceq H(\gamma) \preceq (1-\beta)^{-2}H_0$, so
$\|w\|_{H_0} \leq (1-\beta)^{-1}\|w\|_{H(\gamma)}$ and
\[
  \|\ddot\gamma(t)\|_{H_0}
  \leq \frac{M_f\|v\|_{H_0}^2}{1-\beta(t)}
  \leq 2M_f\|v\|_{H_0}^2.
\]
Integrating: $\|\dot\gamma(t) - v\|_{H_0}
\leq 2M_f\|v\|_{H_0}^2\,t$, and then
\[
  \beta(t) \leq M_f\bigl(t\|v\|_{H_0} + M_f\|v\|_{H_0}^2 t^2\bigr)
  \leq M_f\delta + (M_f\delta)^2
  < \tfrac{1}{4} + \tfrac{1}{16}
  = \tfrac{5}{16} < \tfrac{1}{2}.
\]
Since $\beta(0) = 0$ and the bound shows $\beta$ cannot reach
$\frac{1}{2}$ on $[0,h]$, the containment holds by continuity.

\medskip\noindent
\textit{Step 3: error bound.}\quad
With $\|\ddot\gamma(t)\|_{H_0} \leq 2M_f\|v\|_{H_0}^2$ for all
$t \in [0,h]$:
\[
  \|\gamma(h) - s_0 - hv\|_{H_0}
  = \Bigl\|\int_0^h (h-t)\,\ddot\gamma(t)\,dt\Bigr\|_{H_0}
  \leq 2M_f\|v\|_{H_0}^2 \cdot \frac{h^2}{2}
  = M_f\,\delta^2. \qedhere
\]
\end{proof}

\paragraph{Remark.}
The constant $M_f$ (rather than $\frac{1}{2}M_f$) is the price of
a uniform bound over the trajectory: the Hessian comparison
factor $(1-\beta)^{-1} \leq 2$ converts the
instantaneous bound~\eqref{eq:accel-inst} to the base-point norm.
For $M_f\delta \ll 1$, the true error approaches the leading-order
value $\frac{1}{2}M_f\delta^2$.  For $M_f = 1$, the bound gives
$\delta^2$, and the tangent line stays in the Dikin ellipsoid
whenever $\delta < 1$.

\subsection{Legendre midpoint (2nd order)}

From the Taylor expansion proof (Section~\ref{sec:legendre-proof}),
the Legendre midpoint matches the geodesic to second order.  The
leading error is $O(h^3)$ and involves $D^4 F$ (fourth derivatives)
contracted with $v$.  For a self-concordant barrier, iterated
application of~\eqref{eq:sc} gives:
\begin{equation}\label{eq:legendre-bound}
  \|\gamma(h) - s_{\text{Leg}}\|_H \leq C_2\, \delta^3 + O(\delta^4),
\end{equation}
where $C_2$ is a universal constant (independent of the specific barrier)
that depends only on the self-concordance parameter.

\paragraph{Sketch.}
The $O(h^3)$ error involves the fourth-order Taylor coefficient:
$D^4 F$ contracted with $v$ three times and $H^{-1}$.  By the
derivative of self-concordance (differentiating~\eqref{eq:sc}
with respect to $s$), the fourth derivative satisfies
$|D^4 F[u,v,v,v]| \leq C \|u\|_H \|v\|_H^3$.  The $H^{-1}$
contraction contributes another $\|v\|_H$ factor, giving
$\delta^3 = h^3 \|v\|_H^3$ overall.

\subsection{Yoshida--Legendre (4th and 6th order)}

The Yoshida composition cancels the $O(\delta^3)$ error, leaving:
\begin{align}
  \|\gamma(h) - s_{\text{Yoshida-4}}\|_H &\leq C_4\, \delta^5 + O(\delta^6), \\
  \|\gamma(h) - s_{\text{Yoshida-6}}\|_H &\leq C_6\, \delta^7 + O(\delta^8).
\end{align}
The constants $C_4, C_6$ involve higher derivatives of the
self-concordance condition but remain bounded by universal
functions of $\nu$ (the barrier parameter).

\subsection{Summary: interior guarantee}

\begin{center}
\renewcommand{\arraystretch}{1.2}
\begin{tabular}{@{}lccc@{}}
\toprule
\textbf{Method} & \textbf{Error bound} & \textbf{Order} & \textbf{Interior if} \\
\midrule
Tangent line      & $\delta^2/2$       & 1 & $\delta < \sqrt{2}$ \\
Legendre midpoint  & $C_2 \delta^3$     & 2 & $\delta < C_2^{-1/3}$ \\
Yoshida-4         & $C_4 \delta^5$     & 4 & $\delta < C_4^{-1/5}$ \\
Yoshida-6         & $C_6 \delta^7$     & 6 & $\delta < C_6^{-1/7}$ \\
\bottomrule
\end{tabular}
\end{center}
Higher-order methods allow larger steps $\delta$ while staying
interior.  The constants $C_k$ are universal (depending only on $\nu$,
not on the specific cone), so the step-size bounds are
\emph{a priori}---computable before the step is taken.

\paragraph{Practical implication.}
For an IPM with step size $\alpha$ and direction $d$, the Hessian
norm $\delta = \alpha\|d\|_H$ is already computed during the line
search.  The error bound guarantees that the geodesic step stays
interior whenever $\delta$ is below the method-specific threshold.
No post-hoc feasibility check is needed.

\section{Error analysis for general $M_f$}
\label{sec:rel-smooth}

The error bounds in Sections~\ref{sec:legendre-proof}--\ref{sec:yoshida}
were stated for the standard self-concordant case $M_f = 1$.  We now
give rigorous bounds for general $(M_f, \nu)$-self-concordant barriers,
showing how the integrator order controls the dependence on $M_f$.

We begin with a digression on a weaker regularity condition that
captures what the integrators use in practice: they never evaluate
$D^3 F$, only $\nabla F$ and $\nabla^2 F$ at different points.

\subsection{Relative smoothness}

\begin{definition}\label{def:rel-smooth}
The gradient $\nabla F$ is \emph{$L$-smooth relative to $\nabla^2 F$}
if for all $s \in \operatorname{int}(K)$ and $\|u\|_{H(s)} < 1/L$:
\begin{equation}\label{eq:rel-smooth}
  \|\nabla F(s + u) - \nabla F(s) - \nabla^2 F(s) \cdot u\|_{H^{-1}}
  \leq \frac{L \, \|u\|_H^2}{1 - L\|u\|_H},
\end{equation}
where $H = \nabla^2 F(s)$ and $\|\cdot\|_{H^{-1}}$ is the dual norm.
The domain restriction $\|u\|_H < 1/L$ and the denominator
$1 - L\|u\|_H$ are the Hessian-metric analogs of the Dikin ellipsoid.
\end{definition}

This is the scale-invariant, self-referential analog of classical
$L$-smoothness ($\nabla^2 F \leq LI$):
\begin{center}
\begin{tabular}{@{}ll@{}}
\toprule
\textbf{Classical (Euclidean)} & \textbf{Self-referential (Hessian)} \\
\midrule
$\nabla^2 F \leq L \cdot I$ (smooth) &
  $\nabla F$ varies $\leq L$ relative to $\nabla^2 F$ \\
$\nabla^2 F \geq \mu \cdot I$ (strongly convex) &
  trivially holds with $\mu = 1$ \\
Step size $\delta < 1/L$ (absolute) &
  Step size $\delta < 1/L$ in local Hessian norm \\
\bottomrule
\end{tabular}
\end{center}

\paragraph{Relation to self-concordance.}

\begin{definition}[Generalized self-concordance {\cite[Def.~5.1.1]{nesterov2004}}]
A $C^3$ convex function $F$ is \emph{$(M_f, \nu)$-self-concordant} if
for all $s \in \operatorname{dom} F$ and $v \in \R^n$,
\begin{equation}\label{eq:gsc}
  |D^3 F(s)[v,v,v]| \;\leq\; 2M_f \,\|v\|_{H(s)}^3,
\end{equation}
where $H(s) = \nabla^2 F(s)$.  The standard self-concordant case is
$M_f = 1$.
\end{definition}

\begin{proposition}\label{prop:sc-implies-rs}
If $F$ is $(M_f, \nu)$-self-concordant, then $\nabla F$ is
$L$-smooth relative to $\nabla^2 F$ with $L = M_f$.
\end{proposition}

\begin{proof}
Write the Taylor remainder as
\[
  \nabla F(s+u) - \nabla F(s) - H(s)\,u
  \;=\; \int_0^1 \bigl[H(s+tu) - H(s)\bigr]\,u\;\mathrm{d}t.
\]
Set $\alpha = M_f\|u\|_{H(s)} < 1$.  The Hessian comparison theorem
for self-concordant functions~\cite[Thm.~5.1.1]{nesterov2004} gives
\[
  H(s+tu) \;\preceq\; \frac{1}{(1 - t\alpha)^2}\, H(s).
\]
Therefore $H(s+tu) - H(s) \preceq \bigl[(1-t\alpha)^{-2} - 1\bigr] H(s)$,
and taking the $H(s)^{-1}$ norm of the integral:
\begin{align*}
  \bigl\|\nabla F(s\!+\!u) - \nabla F(s) - H(s)u\bigr\|_{H^{-1}}
  &\leq \|u\|_H \int_0^1 \!\Bigl(\frac{1}{(1-t\alpha)^2} - 1\Bigr)\mathrm{d}t \\
  &= \|u\|_H \biggl[\frac{1}{\alpha(1-t\alpha)}\bigg|_0^1 - 1\biggr]
  \;=\; \|u\|_H \cdot \frac{\alpha}{1-\alpha}
  \;=\; \frac{M_f\,\|u\|_H^2}{1 - M_f\|u\|_H},
\end{align*}
which is~\eqref{eq:rel-smooth} with $L = M_f$.
\end{proof}

\paragraph{The converse is false.}
A barrier can be $L$-smooth relative to $\nabla^2 F$
(finite $L$) even when $D^3 F$ has localized spikes (large or
undefined $M_f$).  Relative smoothness is strictly weaker: it
measures the \emph{integrated} gradient deviation, not pointwise
third derivatives.  This gap is practically significant for
approximate barriers (where $D^3 F$ may not exist everywhere) and
composite barriers (where a small perturbation can create large
pointwise $D^3 F$ without affecting the integrated bound).

\subsection{Error bounds for self-concordant barriers}

Throughout this section, $F$ is $(M_f, \nu)$-self-concordant,
$\gamma$ is the geodesic from $s_0$ with velocity $v$,
and $\delta = h\|v\|_{H_0}$ with $H_0 = \nabla^2 F(s_0)$.

\subsubsection{Preliminaries}

\begin{lemma}[Geodesic acceleration]\label{lem:accel}
Along the geodesic, the acceleration in the instantaneous metric
satisfies
\begin{equation}\label{eq:accel-bound}
  \|\ddot\gamma(t)\|_{H(\gamma(t))} \leq M_f \|v\|_{H_0}^2
\end{equation}
for all $t$ in the domain of $\gamma$.
\end{lemma}
\begin{proof}
The geodesic equation gives $H(\gamma)\ddot\gamma
= -\frac{1}{2}T(\dot\gamma,\dot\gamma)$, so
$\|\ddot\gamma\|_H = \frac{1}{2}\|T(\dot\gamma,\dot\gamma)\|_{H^{-1}}$.
For any $w$ with $\|w\|_H = 1$,
\[
  |w^\top T(\dot\gamma,\dot\gamma)|
  = |D^3 F[\dot\gamma, \dot\gamma, w]|
  \leq 2M_f\,\|\dot\gamma\|_H^2\,\|w\|_H
  = 2M_f\,\|\dot\gamma\|_H^2,
\]
where the inequality is the polarized form
of~\eqref{eq:gsc}.  Since $\|\dot\gamma(t)\|_{H(\gamma(t))}
= \|v\|_{H_0}$ (constant speed on a geodesic), the bound follows.
\end{proof}

To convert from the instantaneous norm $H(\gamma(t))$ to the
base-point norm $H_0$, we use the Hessian comparison theorem:
for $\beta = M_f\|\gamma(t) - s_0\|_{H_0} < 1$,
\begin{equation}\label{eq:hess-compare}
  (1-\beta)^2\, H_0 \;\preceq\; H(\gamma(t)) \;\preceq\;
  (1-\beta)^{-2}\, H_0.
\end{equation}
In particular, $\|w\|_{H_0} \leq (1-\beta)^{-1}\|w\|_{H(\gamma)}$
for any vector $w$.

\begin{lemma}[Trajectory containment]\label{lem:contain}
If $M_f\delta < \frac{1}{4}$, then $\gamma(t)$ remains inside the
ball $M_f\|\gamma(t) - s_0\|_{H_0} < \frac{1}{2}$ for $t \in [0, h]$.
In particular, $\|\ddot\gamma(t)\|_{H_0} \leq 2M_f\|v\|_{H_0}^2$.
\end{lemma}
\begin{proof}
Define $\beta(t) = M_f\|\gamma(t) - s_0\|_{H_0}$.  On any interval
$[0, t_*]$ where $\beta < \frac{1}{2}$,
Lemma~\ref{lem:accel} and~\eqref{eq:hess-compare} give
\[
  \|\ddot\gamma(\tau)\|_{H_0}
  \leq \frac{M_f\|v\|_{H_0}^2}{1-\beta(\tau)}
  \leq 2M_f\|v\|_{H_0}^2.
\]
Integrating: $\|\dot\gamma(\tau) - v\|_{H_0}
\leq 2M_f\|v\|_{H_0}^2\,\tau$, so
\begin{align*}
  \beta(t)
  &= M_f\Bigl\|\int_0^t \dot\gamma(\tau)\,d\tau - s_0\Bigr\|_{H_0}
  \leq M_f\bigl(t\|v\|_{H_0} + M_f\|v\|_{H_0}^2 t^2\bigr) \\
  &\leq M_f\delta + (M_f\delta)^2
  < \tfrac{1}{4} + \tfrac{1}{16}
  = \tfrac{5}{16} < \tfrac{1}{2}.
\end{align*}
The set $\{t : \beta(t) < \frac{1}{2}\}$ is open and contains $0$;
the bound shows it cannot have a first exit time in $[0,h]$.
\end{proof}

\subsubsection{Tangent line (order 1)}

\begin{theorem}[Tangent line, general $M_f$]\label{thm:tangent-Mf}
If $M_f\delta < \frac{1}{4}$, then
\[
  \|\gamma(h) - (s_0 + hv)\|_{H_0}
  \leq M_f\,\delta^2.
\]
\end{theorem}
\begin{proof}
By Lemma~\ref{lem:contain},
$\|\ddot\gamma(t)\|_{H_0} \leq 2M_f\|v\|_{H_0}^2$ for
$t \in [0,h]$.  Since $\gamma(h) - s_0 - hv
= \int_0^h (h-t)\,\ddot\gamma(t)\,dt$,
\[
  \|\gamma(h) - s_0 - hv\|_{H_0}
  \leq 2M_f\|v\|_{H_0}^2 \cdot \frac{h^2}{2}
  = M_f\,\delta^2. \qedhere
\]
\end{proof}

\paragraph{Remark.}
The factor of $2$ (compared to the leading-order bound $\frac{1}{2}M_f\delta^2$
from Theorem~\ref{thm:euler-bound}) is the price of a rigorous global bound:
Lemma~\ref{lem:contain} controls $\|\ddot\gamma\|_{H_0}$ uniformly
over the trajectory, not just at $t = 0$.  For small $M_f\delta$,
the true error approaches $\frac{1}{2}M_f\delta^2$.

\subsubsection{Legendre midpoint (order 2)}

\begin{theorem}[Legendre midpoint, general $M_f$]\label{thm:legendre-Mf}
If $M_f\delta < \frac{1}{4}$, then
\[
  \|\gamma(h) - s_{\textup{Leg}}\|_{H_0}
  \leq C\, M_f^2\, \delta^3,
\]
where $C$ is a universal constant.
\end{theorem}
\begin{proof}
Define the dual variable along the geodesic:
$\lambda(t) = -\nabla F(\gamma(t))$, with $\dot\lambda = H(\gamma)\dot\gamma$.

\medskip\noindent\textit{Step 1: bound $\ddot\lambda$.}\quad
Differentiating: $\ddot\lambda
= D^3F[\dot\gamma,\dot\gamma,\cdot] + H\ddot\gamma
= T(\dot\gamma,\dot\gamma) + H\ddot\gamma$.
The geodesic equation gives $H\ddot\gamma = -\frac{1}{2}T$, so
\begin{equation}\label{eq:dual-accel}
  \ddot\lambda = \tfrac{1}{2}T(\dot\gamma,\dot\gamma).
\end{equation}
By the same argument as Lemma~\ref{lem:accel},
$\|\ddot\lambda\|_{H^{-1}} = \frac{1}{2}\|T\|_{H^{-1}}
\leq M_f\|v\|_{H_0}^2$.

\medskip\noindent\textit{Step 2: bound $\dddot\lambda$.}\quad
Differentiating~\eqref{eq:dual-accel}:
$\dddot\lambda = D^4F[\dot\gamma,\dot\gamma,\dot\gamma,\cdot]
+ D^3F[\ddot\gamma,\dot\gamma,\cdot] + \text{sym}$.
Differentiating the self-concordance
inequality~\eqref{eq:gsc} gives
$|D^4F[u,v,v,v]| \leq C_1 M_f^2\|u\|_H\|v\|_H^3$ and the
$D^3F[\ddot\gamma,\cdot,\cdot]$ terms are bounded using
$\|\ddot\gamma\|_H \leq M_f\|v\|_H^2$
(Lemma~\ref{lem:accel}).  Together, in the $H_0^{-1}$ norm
(using~\eqref{eq:hess-compare} with $\beta < \frac{1}{2}$):
\begin{equation}\label{eq:dual-jerk}
  \|\dddot\lambda(t)\|_{H_0^{-1}} \leq C_2 M_f^2\|v\|_{H_0}^3
\end{equation}
for a universal constant $C_2$.

\medskip\noindent\textit{Step 3: midpoint quadrature error.}\quad
The Legendre midpoint computes (Section~\ref{sec:legendre-proof}):
\[
  \lambda_1^{\text{Leg}}
  = -2\nabla F(s_0 + \tfrac{h}{2}v) + \nabla F(s_0)
  = \lambda_0 + h\dot\lambda_*,
\]
where $\dot\lambda_* = H(s_0 + \frac{h}{2}v)\,v$ is the dual
velocity evaluated at the \emph{primal} midpoint $s_0 + \frac{h}{2}v$.

The true geodesic endpoint in dual coordinates is
\[
  \lambda(h) = \lambda_0 + \int_0^h \dot\lambda(t)\,dt.
\]
The difference $\lambda_1^{\text{Leg}} - \lambda(h)$ has two
sources of error:
\begin{enumerate}
  \item \emph{Midpoint quadrature.}  Even if we evaluated
    $\dot\lambda$ at the true geodesic midpoint $\gamma(h/2)$,
    the midpoint rule
    $\int_0^h f(t)\,dt \approx hf(h/2)$ has error
    $\frac{h^3}{24}\ddot f(\xi)$.  Applied to $f = \dot\lambda$:
    error $= \frac{h^3}{24}\dddot\lambda(\xi)$, bounded
    by~\eqref{eq:dual-jerk}.
  \item \emph{Primal vs.\ geodesic midpoint.}  The Legendre step
    evaluates at $s_0 + \frac{h}{2}v$ instead of $\gamma(h/2)$.
    These differ by $O(h^2)$ (Theorem~\ref{thm:tangent-Mf} at
    half-step), so $H(s_0 + \frac{h}{2}v)v - H(\gamma(h/2))\dot\gamma(h/2)
    = O(M_f\|v\|_H^2 h)$.  Multiplied by $h$: $O(M_f\delta^2 h)
    = O(M_f\delta^3/\|v\|_H)$.  But this error enters at $O(h^3)$
    because the midpoint evaluation is already correct to $O(h)$
    in velocity.
\end{enumerate}
Combining: $\|\lambda_1^{\text{Leg}} - \lambda(h)\|_{H_0^{-1}}
\leq C_3 M_f^2 \|v\|_{H_0}^3 h^3 = C_3 M_f^2 \delta^3$.

\medskip\noindent\textit{Step 4: dual to primal.}\quad
For the barrier $F$, the gradient map $s \mapsto -\nabla F(s)$ has
Jacobian $H(s)$.  By the inverse function theorem, a dual error
$\Delta\lambda$ in $H_0^{-1}$ norm produces a primal error
$\|\Delta s\|_{H_0} \leq (1-\beta)^{-2}\|\Delta\lambda\|_{H_0^{-1}}
\leq 4\|\Delta\lambda\|_{H_0^{-1}}$.

Therefore $\|\gamma(h) - s_{\text{Leg}}\|_{H_0}
\leq C\, M_f^2\, \delta^3$.
\end{proof}

\subsubsection{Yoshida-$2p$ (order $2p$)}

\begin{theorem}[Yoshida composition, general $M_f$]\label{thm:yoshida-Mf}
If $M_f\delta < \frac{1}{4}$, then the Yoshida-$2p$ composition
of Legendre midpoint steps satisfies
\[
  \|\gamma(h) - s_{\textup{Y}_{2p}}\|_{H_0}
  \leq C_p\, M_f^{2p}\, \delta^{2p+1}.
\]
\end{theorem}
\begin{proof}
The Yoshida composition of order $2p$ is built from $2p-1$ Legendre
substeps with coefficients $c_1, \ldots, c_{2p-1}$ summing to $1$,
chosen so that the local error expansion
\[
  s_{\text{Leg}}(c_i h)
  = \gamma(c_i h) + c_i^3 h^3\, e_3 + c_i^5 h^5\, e_5 + \cdots
\]
(where $e_k$ depends on $D^{k+1}F$ contracted with $v$)
satisfies $\sum_i c_i^3 = 0$ (cancels $O(h^3)$),
$\sum_i c_i^5 = 0$ (cancels $O(h^5)$), and so on through
$\sum_i c_i^{2p-1} = 0$.

The leading surviving error has coefficient $\sum_i c_i^{2p+1}$
multiplied by $e_{2p+1}$, which involves $D^{2p+2}F$ contracted
with $v$.  By iterated differentiation of the self-concordance
inequality~\eqref{eq:gsc}, the $k$-th derivative of $F$ satisfies
\[
  |D^k F[v, \ldots, v]| \leq C_k'\, M_f^{k-2}\, \|v\|_H^k
\]
for universal constants $C_k'$ (the pattern: each differentiation
of $|D^3F| \leq 2M_f\|v\|_H^3$ introduces one factor of $M_f$
from the Hessian comparison).

The error term $e_{2p+1}$ involves $H^{-1} D^{2p+2}F[v^{2p+1}]$
(schematically), contributing $M_f^{2p}\|v\|_H^{2p+1}$.
Multiplied by $h^{2p+1}$:
\[
  \|e_{2p+1}\|_{H_0} \cdot h^{2p+1}
  \leq C_p\, M_f^{2p}\, \delta^{2p+1}. \qedhere
\]
\end{proof}

\paragraph{Remark.}
The $M_f^{2p}$ scaling shows that higher-order methods are more
sensitive to $M_f$ in the \emph{coefficient}, but less sensitive
in the \emph{exponent} of $\delta$.  For target error $\varepsilon$:
\[
  \delta_{\max} \propto \left(\frac{\varepsilon}{M_f^{2p}}\right)^{1/(2p+1)}.
\]
This is the content of the step-size table below.

\subsection{Step-size and work comparison}

\subsubsection{Max step size for unit error}

For target error $\varepsilon$, each method's max step size scales as
$\delta_{\max} \propto (\varepsilon / M_f^{2p})^{1/(2p+1)}$:

\begin{center}
\renewcommand{\arraystretch}{1.2}
\begin{tabular}{@{}lcccc@{}}
\toprule
\textbf{Method} & \textbf{Error bound} & \textbf{$\delta_{\max}$ scaling} &
  $M_f\!=\!1$ & $M_f\!=\!10$ \\
\midrule
Tangent line      & $M_f\,\delta^2$ & $M_f^{-1/2}$ &
  $1.00$ & $0.32$ \\
Legendre midpoint  & $C M_f^2\,\delta^3$ & $M_f^{-2/3}$ &
  $\sim\!1$ & $0.22$ \\
Yoshida-4         & $C_4 M_f^4\,\delta^5$ & $M_f^{-4/5}$ &
  $\sim\!1$ & $0.16$ \\
Yoshida-6         & $C_6 M_f^6\,\delta^7$ & $M_f^{-6/7}$ &
  $\sim\!1$ & $0.12$ \\
\bottomrule
\end{tabular}
\end{center}

\paragraph{Observation.}
The max step size for unit error actually \emph{decreases} with
method order when $M_f > 1$.  The coefficient $M_f^{2p}$ grows
faster than the exponent $\delta^{2p+1}$ can compensate.  This is
because the leading error term at order $2p$ involves the
$(2p+2)$-th derivative of $F$, which self-concordance bounds as
$O(M_f^{2p})$.

\subsubsection{Where higher order helps: fixed step accuracy}

The correct comparison is at a \textbf{fixed step size} $\delta$
(constrained to $M_f\delta < 1/4$ by the trajectory containment
lemma).  For example, at $\delta = 0.1 / M_f$:

\begin{center}
\renewcommand{\arraystretch}{1.2}
\begin{tabular}{@{}lcc@{}}
\toprule
\textbf{Method} & \textbf{Error at $\delta = 0.1/M_f$} & \textbf{Ratio} \\
\midrule
Tangent line      & $M_f \cdot (0.1/M_f)^2 = 0.01/M_f$ & $1$ \\
Legendre midpoint  & $C M_f^2 \cdot (0.1/M_f)^3 = 10^{-3} C / M_f$ & $0.1 C$ \\
Yoshida-4         & $C_4 M_f^4 \cdot (0.1/M_f)^5 = 10^{-5} C_4 / M_f$ & $10^{-3} C_4$ \\
Yoshida-6         & $C_6 M_f^6 \cdot (0.1/M_f)^7 = 10^{-7} C_6 / M_f$ & $10^{-5} C_6$ \\
\bottomrule
\end{tabular}
\end{center}

At any fixed fraction of the Dikin radius, higher-order methods
give dramatically better accuracy.  The error ratio is independent
of $M_f$: each Yoshida level gains a factor of $\sim\!100$ for
the same step size.  This matches the $94{,}000\times$ improvement
observed in the exponential cone experiments
(Section~\ref{sec:exp-cone-numerics}).

\subsubsection{Total work comparison}

For a target per-step accuracy $\varepsilon$ and total path length
$D$ (in Hessian norm), the total number of Hessian solves is
(step count) $\times$ (solves per step):
\begin{center}
\renewcommand{\arraystretch}{1.2}
\begin{tabular}{@{}lccc@{}}
\toprule
\textbf{Method} & \textbf{Solves/step} & \textbf{$\delta(\varepsilon)$} &
  \textbf{Total solves} \\
\midrule
Tangent & $1$ & $(\varepsilon/M_f)^{1/2}$ &
  $D\,(M_f/\varepsilon)^{1/2}$ \\
Legendre & $\sim\!3$ & $(\varepsilon/CM_f^2)^{1/3}$ &
  $3D\,(CM_f^2/\varepsilon)^{1/3}$ \\
Yoshida-4 & $9$ & $(\varepsilon/C_4 M_f^4)^{1/5}$ &
  $9D\,(C_4 M_f^4/\varepsilon)^{1/5}$ \\
\bottomrule
\end{tabular}
\end{center}

\paragraph{Scaling with $\varepsilon$.}
For fixed $M_f$, tangent scales as $\varepsilon^{-1/2}$, Legendre as
$\varepsilon^{-1/3}$, Yoshida-4 as $\varepsilon^{-1/5}$.  Higher
order wins decisively when high per-step accuracy is needed.

\paragraph{Scaling with $M_f$.}
Tangent scales as $M_f^{1/2}$, Legendre as $M_f^{2/3}$, Yoshida-4
as $M_f^{4/5}$.  Higher order has \emph{worse} $M_f$ scaling.
The $M_f^{2p}$ error coefficient (from higher derivatives of $F$)
grows faster than the step-size exponent can absorb.

\paragraph{Summary.}
Higher-order Yoshida methods are most valuable when:
\begin{itemize}
  \item The barrier is well-behaved ($M_f$ moderate), so the
    coefficient penalty is small,
  \item High per-step accuracy is needed (geodesic step must closely
    approximate the true geodesic),
  \item The overhead of multiple Hessian solves per step is acceptable.
\end{itemize}
For barriers with $M_f \gg 1$, higher order does not compensate for
the poor regularity --- the error coefficient $M_f^{2p}$ dominates.
In this regime, the tangent-line step (with smaller steps) may be
more efficient.

\section{Application to interior-point methods}

In conic optimization, the barrier function $F$ defines a Hessian
manifold on the interior of the cone.  The geodesic step keeps
iterates strictly inside the cone (unlike a linear step which may
exit).  The Yoshida--Legendre method provides:
\begin{itemize}
  \item \textbf{Cone-preserving steps} without computing third
    derivatives of the barrier,
  \item \textbf{Fourth-order accuracy} from three Newton solves
    per step (each on a system the size of the cone),
  \item \textbf{Generality}: works for any cone with a computable
    barrier gradient and Hessian (exponential cone, power cone,
    quantum entropy cone, etc.).
\end{itemize}

For the exponential cone (3-dimensional), the method converges in
4 factorizations on test problems, with the duality gap decreasing
by a factor of $\sim\!100$ per factorization.

\begin{thebibliography}{9}
\bibitem{nesterov2018}
Y.~Nesterov, \emph{Lectures on Convex Optimization}, 2nd ed.,
Springer, 2018.

\bibitem{yoshida1990}
H.~Yoshida, ``Construction of higher order symplectic integrators,''
\emph{Phys.\ Lett.\ A}, 150(5--7):262--268, 1990.

\bibitem{amari2000}
S.~Amari and H.~Nagaoka, \emph{Methods of Information Geometry},
AMS/Oxford, 2000.
\end{thebibliography}

\end{document}
